% Chapter 5, Topic _Linear Algebra_ Jim Hefferon
%  http://joshua.smcvt.edu/linearalgebra
%  2016-Jun-17
% \documentclass{article}
% \usepackage{graphicx}
% \usepackage{amsmath}

% \newcommand{\definend}[1]{\textit{#1}}
\topic{Osilator Terkopel}
\index{Osilator terkopel|(}
% \begin{document}

Ini adalah sebuah \definend{bandul Wilberforce}.\index{bandul Wilberforce}
Pada pegasnya tergantung sebuah massa, atau bandul.
Dorong bandul itu sedikit ke atas lalu lepaskan; bandul akan berosilasi naik turun.
\begin{center}
  \includegraphics{jc/asy/wilber.pdf}
\end{center}
Namun, jika perangkat disetel dengan tepat, sesuatu yang menakjubkan terjadi.
Setelah beberapa detik, selain bergerak naik
turun, massa mulai berputar
mengelilingi sumbu yang melalui tengah pegas.
Putaran ini membesar sampai geraknya  
hampir seluruhnya berupa rotasi, dengan sangat sedikit gerak naik turun.
Mungkin lima detik kemudian, geraknya kembali berkembang menjadi gabungan keduanya.
Setelah beberapa waktu lagi, keteraturan muncul kembali.
Yang mengagumkan, kini
geraknya hampir seluruhnya vertikal.
Hal ini terus berulang: perangkat berganti-ganti antara
selang gerak vertikal murni dan selang gerak rotasi murni,
dengan gerak campuran di antaranya. 
(Carilah ``wilberforce pendulum video'' di internet untuk melihat 
beberapa peragaan yang sangat bagus.) 

Setiap keadaan gerak murni merupakan \definend{modus normal} osilasi.
Kita akan menganalisis perilaku perangkat ini ketika berada dalam modus normal. 
Semuanya berkaitan dengan nilai eigen.

\begin{center}
  \includegraphics{jc/asy/wilber001.pdf}
\end{center}

Tuliskan~$x(t)$ untuk gerak vertikal terhadap waktu dan 
$\theta(t)$ untuk gerak rotasinya.
Tetapkan sistem koordinat sehingga pada posisi diam $x=0$ dan~$\theta=0$, 
nilai $x$ positif mengarah ke atas, dan nilai $\theta$ positif
menunjukkan arah berlawanan jarum jam bila dilihat dari atas.

Kita mulai dengan memodelkan gerak sebuah massa pada 
pegas yang dibatasi agar tidak terpuntir.
Kasus ini lebih sederhana karena hanya ada satu gerak, yakni satu derajat kebebasan.
Letakkan massa pada posisi diam, lalu 
dorong ke atas untuk memampatkan pegas.
Hukum Hooke menyatakan bahwa untuk simpangan kecil,
gaya pemulih 
sebanding dengan simpangan tersebut, $F=-k\cdot x$.
Konstanta~$k$ adalah
\definend{kekakuan} pegas.

Hukum Newton menyatakan bahwa gaya sebanding dengan 
percepatan yang ditimbulkannya, $F=m\cdot d^2\,x(t)/dt^2$.
Konstanta kesebandingan $m$ adalah \definend{massa}
benda.
Menggabungkan Hukum Hooke dan Hukum Newton 
menghasilkan persamaan diferensial yang menyatakan 
gerak massa, $m\cdot d^2\,x(t)/dt^2=-k\cdot x(t)$.
Kita memilih bentuk dengan semua variabel pada satu ruas.
\begin{equation*}
  m\cdot \frac{d^2\,x(t)}{dt^2}+k\cdot x(t)=0
  \tag{$*$}
\end{equation*}

Intuisi fisis kita mengatakan bahwa bandul berosilasi seiring waktu.
Bandul mulai dari posisi tinggi, sehingga 
grafiknya akan tampak seperti ini.
\begin{center}
  \includegraphics{jc/asy/wilber003.pdf}
\end{center}
Tentu saja, bentuk ini menyerupai grafik kosinus, dan kita melihat 
bahwa persamaan diferensial gerak~($*$) dipenuhi oleh
$x(t)=\cos \omega t$, dengan $\omega=\sqrt{k/m}$,
karena $dx/dt=-\omega\cdot \sin \omega t$
dan $d^2x/dt^2=-\omega^2\cdot \cos \omega t$.
Di sini, $\omega$ adalah \definend{frekuensi sudut}.
Besaran ini menentukan periode osilasi: jika $\omega=1$, maka
periodenya~$2\pi$, sedangkan jika $\omega=2$, periodenya~$\pi$, dan seterusnya.

Kita dapat memberikan solusi yang lebih umum untuk~($*$).
Untuk amplitudo umum, kita menambahkan faktor~$A$ di depan,
$x(t)=A\cdot \cos \omega t$.
Kita juga dapat memasukkan pergeseran fase, 
sehingga pencatatan waktu tidak harus dimulai ketika
massa berada di posisi tinggi, yaitu 
$x(t)=A\cos (\omega t+\phi)$.
Inilah persamaan \definend{gerak harmonik sederhana}.

Sekarang kita kembali meninjau pasangan gerak yang terkopel, 
yaitu gerak vertikal dan rotasi.
Keduanya saling berinteraksi karena pegas yang terpuntir akan 
sedikit memanjang atau memendek;
misalnya, pengaruhnya dapat berlangsung seperti pada gambar ini atau sebaliknya,
bergantung pada pegasnya.  
\begin{center}
  \includegraphics{jc/asy/wilber002.pdf}
\end{center}
Selain itu, pegas yang diregangkan atau dimampatkan dari posisi
diamnya akan sedikit terpuntir.
Interaksi keduanya menghasilkan \definend{osilasi terkopel}.

Untuk memahami bagaimana interaksi ini dapat menghasilkan perilaku dramatis
yang kita lihat dalam modus normal,
bayangkan massa berputar ke arah yang 
membuat pegas memanjang.
Jika pada saat yang sama gerak vertikal membuat pegas 
memendek, 
maka penumpangtindihan kedua gerak tersebut
dapat membuat keduanya hampir saling meniadakan.
Akibatnya, bandul hampir tidak bergerak secara vertikal dan hanya terpuntir.
Pada perangkat yang disetel dengan tepat, keadaan ini dapat berlangsung selama beberapa detik.

``Disetel dengan tepat'' berarti bahwa  
periode gerak vertikal murni sama dengan, 
atau mendekati, periode gerak rotasi murni.
Dengan demikian, peniadaan tersebut akan berlangsung selama beberapa waktu. 

Interaksi antargerak juga dapat menghasilkan perilaku modus normal yang lain, 
yakni bandul bergerak terutama secara vertikal tanpa banyak berotasi,
jika gerak pegas~$x(t)$ 
menghasilkan puntiran yang melawan puntiran bandul~$\theta(t)$.
Bandul hampir berhenti berotasi, sehingga
geraknya hampir seluruhnya vertikal.

Untuk memperoleh persamaan gerak pada kasus dua derajat kebebasan ini,
kita membuat asumsi yang sama seperti pada kasus satu derajat kebebasan,
yaitu bahwa untuk simpangan kecil, gaya pemulih
sebanding dengan simpangannya. 
Namun, kini kita menerapkan asumsi itu pada
gerak vertikal maupun gerak rotasi.
Misalkan konstanta kesebandingan untuk gerak rotasi adalah~$\kappa$.
Serupa dengan itu, kita juga menggunakan Hukum Newton bahwa gaya sebanding dengan 
percepatan untuk gerak rotasi, dan mengambil
konstanta kesebandingannya sebagai~$I$. 

Yang paling penting, kita menambahkan kopling antara kedua
gerak, yang kita anggap sebanding dengan masing-masing gerak dengan konstanta 
kesebandingan~$\epsilon/2$.

Dengan demikian, kita memperoleh sistem dua persamaan diferensial:
persamaan pertama untuk gerak vertikal dan persamaan kedua
untuk rotasi.
Persamaan-persamaan ini mendeskripsikan perilaku sistem terkopel pada setiap waktu~$t$.
\begin{equation*}
  \begin{split}
  m\cdot\frac{d^2\,x(t)}{dt^2}
      +k\cdot x(t)+\frac{\epsilon}{2}\cdot \theta(t) 
      &=0  \\
  I\cdot\frac{d^2\,\theta(t)}{dt^2}
      +\kappa\cdot \theta(t)+\frac{\epsilon}{2}\cdot x(t)  
      &=0
  \end{split}
  \tag{$**$}
\end{equation*}
Kita akan menggunakannya untuk menganalisis
perilaku sistem ketika berada dalam modus normal.

Pertama-tama, tinjau gerak-gerak tak terkopel, yaitu gerak yang diberikan oleh
persamaan tanpa suku~$\epsilon$.
Tanpa suku tersebut, persamaan ini mendeskripsikan fungsi harmonik sederhana;
kita menuliskan~$\omega_x^2$ untuk~$k/m$ dan $\omega_\theta^2$ untuk $\kappa/I$.
Di atas kita telah berargumen bahwa untuk mengamati perilaku hampir diam tersebut,
perangkat harus disetel agar periodenya sama, $\omega_x^2=\omega_\theta^2$. 
Tuliskan~$\omega_0^2$ untuk nilai bersama itu.

Sekarang tinjau gerak terkopel $x(t)$ dan~$\theta(t)$.
Berdasarkan prinsip yang sama, untuk mengamati perilaku hampir diam tersebut, kita ingin keduanya 
selaras; misalnya, rotasi mencapai
puncaknya ketika regangan mencapai puncaknya.
Artinya, dalam modus normal, kedua osilasi mempunyai
frekuensi sudut~$\omega$ yang sama.
Mengenai pergeseran fase, seperti yang juga telah kita bahas, ada dua
kasus: puntiran yang ditimbulkan gerak pegas searah
dengan puntiran akibat osilasi rotasi, atau kedua puntiran itu berlawanan arah.
Dalam kedua kasus, agar diperoleh modus normal, puncak-puncaknya harus berimpit.
\begin{align*}
  x(t) &= A_1\cos(\omega t+\phi)  
       &x(t) &= A_1\cos(\omega t+\phi)   \\  
  \theta(t) &= A_2\cos(\omega t+\phi)  
       &\theta(t) &= A_2\cos(\omega t+(\phi+\pi))    
\end{align*}
Kita akan mengerjakan kasus di sebelah kiri dan menyerahkan kasus lainnya sebagai latihan.

Kita ingin mencari nilai~$\omega$ yang mungkin.
Hitung turunan kedua masing-masing: 
\begin{equation*}
  \frac{d^2\,x(t)}{dt^2}=-A_1\omega^2\cos(\omega t+\phi)
  \qquad
  \frac{d^2\,\theta(t)}{dt^2}=-A_2\omega^2\cos(\omega t+\phi)
\end{equation*}
dan substitusikan ke persamaan gerak~($**$).
\begin{align*}
  m\cdot(-A_1\omega^2\cos(\omega t+\phi))
      +k\cdot (A_1\cos(\omega t+\phi))
      +\frac{\epsilon}{2}\cdot (A_2\cos(\omega t+\phi)) 
      &=0  \\
  I\cdot(-A_2\omega^2\cos(\omega t+\phi))
      +\kappa\cdot (A_2\cos(\omega t+\phi))
      +\frac{\epsilon}{2}\cdot (A_1\cos(\omega t+\phi))  
      &=0  
\end{align*}
Keluarkan faktor $\cos(\omega t+\phi)$, lalu bagi persamaan pertama dengan~$m$
dan persamaan kedua dengan~$I$.
\begin{align*}
  \big(\frac{k}{m}-\omega^2\big)\cdot A_1 
      +\frac{\epsilon}{2m}\cdot A_2 
      &=0  \\
  \big(\frac{\kappa}{I}-\omega^2\big)\cdot A_2
      +\frac{\epsilon}{2I}\cdot A_1  
      &=0  
\end{align*}
Kita mengasumsikan bahwa $k/m=\omega_{x}^2$ dan
$\kappa/I=\omega_{\theta}^2$ sama, serta menuliskan $\omega_0^2$
untuk nilai bersama itu.
Lakukan substitusi tersebut, lalu nyatakan kembali sebagai persamaan matriks.
\begin{equation*}
  \begin{mat}
    \omega_0^2-\omega^2    &\epsilon/2m  \\
    \epsilon/2I &\omega_0^2-\omega^2 
  \end{mat}
  \colvec{
    A_1 \\
    A_2}
  =
  \begin{mat}
    0 \\
    0
  \end{mat}
\end{equation*}
Jelas bahwa sistem ini memiliki solusi trivial~$A_1=0$, $A_2=0$,
yang bersesuaian dengan keadaan ketika massa diam.
Kita ingin mengetahui frekuensi~$\omega$ yang membuat sistem ini memiliki
solusi nontrivial. 
\begin{equation*}
  \begin{mat}
    \omega_0^2    &\epsilon/2m  \\
    \epsilon/2I &\omega_0^2
  \end{mat}
  \colvec{
    A_1 \\
    A_2}
  =
  \omega^2
  \begin{mat}
    A_1 \\
    A_2
  \end{mat}
\end{equation*}
Kuadrat frekuensi sudut modus normal, yakni~$\omega^2$, merupakan nilai eigen matriks tersebut.

Untuk menghitungnya,
ambil determinannya dan samakan dengan nol.
\begin{equation*}
  \begin{vmatrix}
    \omega_0^2-\omega^2    &\epsilon/2m  \\
    \epsilon/2I &\omega_0^2-\omega^2
  \end{vmatrix}=0 \quad\Longrightarrow\quad
  \omega^4-(2\omega_0^2)\,\omega^2 +(\omega_0^4
  -\frac{\epsilon^2}{4mI})=0
\end{equation*}
Persamaan itu berbentuk kuadrat dalam~$\omega^2$.  
Terapkan rumus penyelesaian persamaan kuadrat, $(-b\pm\sqrt{b^2-4ac})/(2a)$.
\begin{equation*}   % \rule[-.3\baselineskip]{0pt}{\baselineskip}
  \omega^2 
  =\frac{2\omega_0^2
    \pm\sqrt{\smash[b]{\strut} 4\omega_0^4
      -4(\omega_0^4-\epsilon^2/4mI)}}{2}
  =\omega_0^2\pm\frac{\epsilon}{2\sqrt{mI}}
\end{equation*}
Nilai $\epsilon/\sqrt{mI}$ sering dituliskan sebagai
$\omega_B$, sehingga
$
\omega^2 =\omega_0^2\pm\omega_B/2
$.
Besaran ini adalah \definend{frekuensi layangan}; nilainya sama dengan selisih
antara kuadrat kedua frekuensi modus normal.

Walaupun penjelasannya berada di luar cakupan kita, rumus umum untuk 
gerak bandul merupakan kombinasi linear dari gerak-gerak pada 
modus normal.
Jadi, gerak bandul sepenuhnya ditentukan oleh nilai-nilai eigen
matriks di atas.
Lihat \cite{BergMarshall}.

\begin{exercises}
\item Gunakan rumus kosinus jumlah sudut untuk memperoleh rumus gerak
  harmonik sederhana yang lebih umum lagi.
  \begin{answer}
    Rumus jumlah sudut untuk fungsi kosinus adalah
    $\cos(\alpha+\beta)=\cos(\alpha)\cos(\beta)-\sin(\alpha)\sin(\beta)$.
    Uraikan $A\cos(\omega t+\phi)$ menjadi $A\cdot[\cos(\omega
    t)\cos(\phi)-\sin(\omega t)\sin(\phi)]$.  Karena $\cos(\phi)$ dan
    $\sin(\phi)$ tidak berubah terhadap~$t$, kita memperoleh solusi umum
    $x(t)=B\cos(\omega t)+C\sin(\omega t)$.
  \end{answer}
\item Carilah vektor-vektor eigen yang terkait dengan nilai-nilai eigen tersebut.
  \begin{answer}
    Kita akan mengerjakan akar pertama; akar kedua dikerjakan dengan cara serupa.  
    Kita mempunyai persamaan berikut.
    \begin{equation*}
      \begin{mat}
        \omega_0^2-\omega^2  &\epsilon/2m \\
        \epsilon/2I         &\omega_0^2-\omega^2
      \end{mat}
      \colvec{A_1 \\ A_2}
      =\colvec{0 \\ 0}
    \end{equation*}
    Substitusikan akar pertama~$\omega^2=\omega_0^2+\epsilon/(2\sqrt{mI})$.
    Kedua persamaan itu memberikan syarat yang sama, jadi kita cukup meninjau persamaan pertama.
    \begin{equation*}
      -\frac{\epsilon}{2\sqrt{mI}}\cdot A_1+\frac{\epsilon}{2m}\cdot A_2=0
    \end{equation*}
    Ini adalah garis yang melalui titik asal pada bidang, sehingga untuk menentukannya kita
    hanya perlu mencari rasio antara variabel pertama dan kedua.
    \begin{equation*}
      \frac{A_2}{A_1}=\frac{\epsilon/(2\sqrt{mI})}{\epsilon/(2m)}
                     =\sqrt{\frac{m}{I}}
    \end{equation*}
    Jadi, inilah ruang eigen yang terkait dengan nilai eigen pertama.
    \begin{equation*}
      \set{\colvec{A_1 \\ \sqrt{m/I}\cdot A_1}\suchthat A_1\in\C}
    \end{equation*}
  \end{answer}
\item Carilah nilai-nilai $\omega$ untuk kasus ketika 
  $x(t)=A_1\cos(\omega t+\phi)$
  dan $\theta(t)=A_2\cos(\omega t+(\phi+\pi))$.
  \begin{answer}
    Perhatikan bahwa $\cos(\omega t+(\phi+\pi))=-\cos(\omega t+\phi)$.
    Hitung turunan untuk~$x(t)$: 
    \begin{equation*}
      \frac{d\,x(t)}{dt}=-A_1\omega\cdot\sin(\omega t+\phi)
      \qquad
      \frac{d^2\,x(t)}{dt^2}=-A_1\omega^2\cdot\cos(\omega t+\phi)
    \end{equation*}
    dan untuk~$\theta(t)$: 
    \begin{equation*}
      \frac{d\,\theta(t)}{dt}=A_2\omega\cdot\sin(\omega t+\phi)
      \qquad
      \frac{d^2\,\theta(t)}{dt^2}=A_2\omega^2\cdot\cos(\omega t+\phi)
    \end{equation*}
    lalu substitusikan ke persamaan gerak~($**$).
    \begin{align*}
      mA_1\cdot(-\omega^2\cos(\omega t+\phi))
          +kA_1\cdot (\cos(\omega t+\phi))
          +\frac{\epsilon}{2}A_2\cdot (-\cos(\omega t+\phi)) 
          &=0  \\
      IA_2\cdot(\omega^2\cos(\omega t+\phi))
          +\kappa A_2\cdot (-\cos(\omega t+\phi))
          +\frac{\epsilon}{2}A_1\cdot (\cos(\omega t+\phi))  
          &=0  
    \end{align*}
    Keluarkan faktor $\cos(\omega t+\phi)$: 
    \begin{align*}
     \big(m(-\omega^2)\big)\cdot A_1
         +kA_1 
         -\frac{\epsilon}{2}\cdot A_2 
         &=0  \\
     \big(I\omega^2\big)\cdot A_2
         -\kappa A_2
         +\frac{\epsilon}{2}\cdot A_1  
         &=0  
   \end{align*}
    lalu bagi persamaan pertama dengan~$m$ dan persamaan kedua dengan~$I$.
    \begin{align*}
     \big(\frac{k}{m}-\omega^2\big)\cdot A_1 
         -\frac{\epsilon}{2m}\cdot A_2 
         &=0  \\
     \frac{\epsilon}{2I}\cdot A_1  
     -\big(\frac{\kappa}{I}-\omega^2\big)\cdot A_2
         &=0  
   \end{align*}
   Kita mengasumsikan bahwa $k/m=\omega_{x}^2$ dan
   $\kappa/I=\omega_{\theta}^2$ sama, serta menuliskan $\omega_0^2$
   untuk nilai bersama itu.
   Lakukan substitusi tersebut, lalu nyatakan kembali sebagai persamaan matriks.
   \begin{equation*}
     \begin{mat}
       \omega_0^2-\omega^2    &-\epsilon/2m  \\
       \epsilon/2I &\omega^2-\omega_0^2 
     \end{mat}
     \colvec{
       A_1 \\
       A_2}
     =
     \colvec{
       0 \\
       0}
   \end{equation*}
   Kita mencari frekuensi~$\omega$ yang membuat sistem ini memiliki
   solusi nontrivial. 
   \begin{equation*}
      \begin{vmatrix}
       \omega_0^2-\omega^2    &-\epsilon/2m  \\
       \epsilon/2I           &\omega^2-\omega_0^2 
      \end{vmatrix}
      =0
    \end{equation*}
    Determinan ini sama dengan yang kita hitung dalam uraian utama topik ini,
    kecuali bahwa kolom keduanya dikalikan dengan~$-1$.
    Mengalikan sebuah baris atau kolom dengan suatu skalar akan mengalikan seluruh determinan
    dengan skalar tersebut.
    Jadi, determinan ini adalah negatif dari determinan dalam uraian utama topik ini.
    Namun, karena determinannya disamakan dengan nol, perbedaan itu tidak berpengaruh.
    Akar-akarnya sama dengan yang diperoleh dalam uraian utama topik ini. 
  \end{answer}
\item Buatlah bandul Wilberforce dari Slinky Jr dan sebuah kaleng sup.
  Anda dapat mengebor dua atau empat lubang pada kaleng untuk memasang baut,
  yang dapat digunakan untuk menyesuaikan
  momen inersia kaleng agar periode gerak vertikal dan 
  gerak rotasi berimpit.
  \begin{answer}
    Lihat \cite{Mewes}.
  \end{answer}
\end{exercises}
\index{Osilator terkopel|)}
\endinput
% \end{document}
